People load the same names and pages in patterns. This paper predicts each person's next DNS lookup on macOS and next same-site page in Chrome from their own history, and does the waiting early when the expected time saved exceeds the measured cost. One predictor, a Bayesian mixture of order-0, order-1, order-2 and hour-of-day models with Krichevsky--Trofimov estimates and a daily-tuned decay, runs identically in Python and JavaScript. In offline replay of \V{r0.ntraces} traces it had lower log-loss than an order-1 Markov baseline on all \V{r0.wins_markov_bits}, and lower than PPM-C on \V{r0.wins_ppmc_bits}. On a campus network, API calls in later rounds were \V{r3.warm.paired}\,ms faster per call [\V{r3.warm.paired.lo}, \V{r3.warm.paired.hi}], first calls were \V{r3.cold.paired}\,ms slower, and the shared campus resolver beat a local resolver on median lookup time. The Chrome extension holds only the \texttt{storage} permission, stores only hashes, and never preloads action links.
